\section*{第 \ref{ch:g3:solids-liquids-gases} 章 --- 固体、液体、气体}

\begin{solution}{exo:g3:solids-liquids-gases:1}
固体：硬币、木珠。液体：橙汁、油。
气体：轮胎里的空气，透明缝隙里看不见的水蒸气。
\end{solution}

\begin{solution}{exo:g3:solids-liquids-gases:2}
固体有自己的形状，液体则没有。从玻璃杯倒入碗中，
液体的\emph{量}保持不变。
\end{solution}

\begin{solution}{exo:g3:solids-liquids-gases:3}
量相同——只是形状变了，从瓶子形状变成
宽碗的形状。
\end{solution}

\begin{solution}{exo:g3:solids-liquids-gases:4}
两幅图中，水面都平整、水平，与地面平行——
在倾斜的水槽中，水面并不随杯子倾斜；
水聚集在低的一侧。
\end{solution}

\begin{solution}{exo:g3:solids-liquids-gases:5}
小窗口里的液体总是保持水平，所以气泡
能指示“向上”的方向：气泡位于两条刻线之间，表示搁板
与水平的液面相符。这里用到的定律是：静止的液体
保持水平。
\end{solution}

\begin{solution}{exo:g3:solids-liquids-gases:6}
气体家族。煎饼的气味属于气体，气体会扩散，
充满所给的整个空间——厨房、楼梯间、卧室。
\end{solution}

\begin{solution}{exo:g3:solids-liquids-gases:7}
沙子由许多微小的固体颗粒组成。这群颗粒会流动，并取
所堆积的形状，但每一颗颗粒都完整地保持自己的形状——
因此每一颗颗粒都是真正的固体。
\end{solution}

\begin{solution}{exo:g3:solids-liquids-gases:8}
蜂蜜会（缓慢地）流动，取容器的形状，并保持
自己的量：蜂蜜在这些问题上像液体一样——只是一种懒惰的液体。
\end{solution}

\begin{solution}{exo:g3:solids-liquids-gases:9}
保持自己的形状：固体（勺子）。取容器的形状，但
保持自己的量：液体（牛奶）。充满所有空间：气体（厨房里的
空气）。
\end{solution}

\begin{solution}{exo:g3:solids-liquids-gases:10}
冰熔化时，物质不再保持自己的\emph{形状}（但
仍保持自己的量）。水蒸发时，它不再保持
甚至不再保持自己的\emph{所在位置}：作为气体，它会扩散到
所有可以到达的空间里。
\end{solution}

\begin{solution}{exo:g3:solids-liquids-gases:11}
橡胶薄膜决定了形状：它是一个被拉伸的固体笼子，
被压缩的空气充满笼子所取形状的每个角落。
戳破气球后，空气立即扩散到整个房间里——
这说明空气本来就没有自己的形状。
\end{solution}
